\begin{exercisechunk}
\item \textbf{Parity requires exponentially many DNF terms.}\suppmark
\label[exercise]{ex:l14-parity-dnf}
\exerciseprofile{2}{1}{2}{\exproof\quad\exconstruction\quad\exinterpretation}{%
Prove that changing the allowed representation can change parity from a
linear-size computation to an exponential-size formula.}
The parity function has a linear-size XOR circuit. This exercise proves that
the same function has no comparably small disjunctive normal form, so the
representation language is part of every size comparison.

Fix an integer \(r\geq 1\), and let
\[
 \operatorname{PARITY}_r(x)
 =x_1\oplus\cdots\oplus x_r,
 \qquad x\in\{0,1\}^r.
\]
A \emph{term} is a conjunction of literals, where a literal is \(x_i\) or
\(\operatorname{NOT}(x_i)\). A \emph{DNF} is a disjunction of terms. Terms
containing both \(x_i\) and \(\operatorname{NOT}(x_i)\) may be removed because
they are never satisfied.
\begin{enumerate}[(a),beginpenalty=10000]
\item Prove that exactly \(2^{r-1}\) inputs have parity one.
\item Let a satisfiable term occur in a DNF that computes
\(\operatorname{PARITY}_r\). Prove that the term must contain a literal for
every input coordinate. Deduce that the term is satisfied by at most one
input.
\item Prove that every DNF for \(\operatorname{PARITY}_r\) has at least
\(2^{r-1}\) terms. Construct a DNF with exactly \(2^{r-1}\) terms, and
therefore prove that the minimum number of terms is exactly \(2^{r-1}\).
\item Explain why the lower bound concerns the DNF representation and does not
contradict the existence of a linear-size XOR circuit.
\end{enumerate}
\end{exercisechunk}
